% GENERATED FILE. Edit canonical lessons, then run npm run generate:books.
% canonical-source-hash: fnv1a64:eaad61ae04cb7b3f
% canonical-lessons: JA-C512-henka-suru, JA-C512-kirikaeru, JA-C512-hayamaru, JA-C512-kasanaru, JA-C512-higawari

\chapter{When Things Change}
\label{ch:ja-when-things-change}
\hlchaptermodality{512}

\begin{chapteropening}
\textbf{By the end of this chapter:} \emph{I can say what I will do if plans change.}

Complete the last lesson of chapter 512: I can say what I will do if plans change.
\end{chapteropening}

\section[\texorpdfstring{\ja{へんかする} (henka suru)}{へんかする (henka suru)}]{\ja{へんかする} (henka suru) — to change; to vary}

\label{lesson:JA-C512-henka-suru}



\begin{quote}
Before the new one: say the Japanese for how well something works, then the Japanese for the origin.
\end{quote}

\subsection*{You'll want to know: \ja{へんかする}}
\textbf{\ja{へんかする}} — \emph{henka suru} — \textquotedblleft{}to change; to vary\textquotedblright{}.

\begin{quote}
\textbf{\ja{もし} \ja{てんきが} \ja{へんかしたら、しあいは} \ja{ちゅうしでしょう。}} — \emph{moshi tenki ga henkashitara, shiai wa chūshi deshō} — \textquotedblleft{}If the weather changes, the match will probably be called off.\textquotedblright{}
\end{quote}

Polite \textbf{\ja{へんかします}}, \emph{henkashimasu}; a \ja{する} verb. Past (polite): \textbf{\ja{まちが} \ja{へんかしました。}}, \emph{machi ga henkashimashita}.

\begin{cousinweb}
\textbf{\ja{へんかする}} (\emph{henka suru}) is the Sino-Japanese (kango) noun \emph{henka} with the native verb \emph{suru}, \textquotedblleft{}to do\textquotedblright{}: two Chinese characters read with their on'yomi, readings borrowed from Chinese. They are \emph{hen}, \textquotedblleft{}change\textquotedblright{} (Middle Chinese \emph{pjenH}) and \emph{ka}, \textquotedblleft{}change\textquotedblright{} (\emph{xwaeH}) (KANJIDIC2; Middle Chinese in Baxter's transcription).
\end{cousinweb}

\begin{grammarlens}[title={What you've built}]
One more word for saying what happens if.
\end{grammarlens}

\subsection*{Your turn}
\begin{itemize}
\raggedright
  \item \emph{Say it:} \emph{henka suru}
  \item \emph{Say it:} \emph{henka suru}, once more
  \item \emph{Say it:} \emph{moshi tenki ga henkashitara, shiai wa chūshi deshō} — If the weather changes, the match will probably be called off.
  \item \emph{Recall:} say the Japanese for how well something works, then the Japanese for the origin, then say \emph{henka suru} again
\end{itemize}

\begin{tcolorbox}[breakable,colback=teal!4,colframe=teal!35!black,title={Before you move on}]
What does \ja{へんかする} mean? (\textquotedblleft{}To change; to vary\textquotedblright{}.) Say it once more.
\end{tcolorbox}
\section[\texorpdfstring{\ja{きりかえる} (kirikaeru)}{きりかえる (kirikaeru)}]{\ja{きりかえる} (kirikaeru) — to switch (over)}

\label{lesson:JA-C512-kirikaeru}



\begin{quote}
Before the new one: say the Japanese for the origin, then the Japanese for to change.
\end{quote}

\subsection*{You'll want to know: \ja{きりかえる}}
\textbf{\ja{きりかえる}} — \emph{kirikaeru} — \textquotedblleft{}to switch (over)\textquotedblright{}.

\begin{quote}
\textbf{\ja{もし} \ja{さむく} \ja{なったら、ふゆの} \ja{ふくに} \ja{きりかえます。}} — \emph{moshi samuku nattara, fuyu no fuku ni kirikaemasu} — \textquotedblleft{}If it gets cold, I will switch to winter clothes.\textquotedblright{}
\end{quote}

Polite \textbf{\ja{きりかえます}}, \emph{kirikaemasu}; an ichidan verb. Past (polite): \textbf{\ja{みせが} \ja{ふゆの} \ja{ふくに} \ja{きりかえました。}}, \emph{mise ga fuyu no fuku ni kirikaemashita}.

\begin{cousinweb}
\textbf{\ja{きりかえる}} (\emph{kirikaeru}) is a native Japanese word (wago), inherited rather than borrowed from Chinese, built from \emph{kiru}, \textquotedblleft{}to cut\textquotedblright{}, and \emph{kaeru}, \textquotedblleft{}to change\textquotedblright{}; its kanji was attached later as a kun'yomi (JMdict, KANJIDIC2).
\end{cousinweb}

\begin{grammarlens}[title={What you've built}]
One more word for saying what happens if.
\end{grammarlens}

\subsection*{Your turn}
\begin{itemize}
\raggedright
  \item \emph{Say it:} \emph{kirikaeru}
  \item \emph{Say it:} \emph{kirikaeru}, once more
  \item \emph{Say it:} \emph{moshi samuku nattara, fuyu no fuku ni kirikaemasu} — If it gets cold, I will switch to winter clothes.
  \item \emph{Recall:} say the Japanese for the origin, then the Japanese for to change, then say \emph{kirikaeru} again
\end{itemize}

\begin{tcolorbox}[breakable,colback=teal!4,colframe=teal!35!black,title={Before you move on}]
What does \ja{きりかえる} mean? (\textquotedblleft{}To switch (over)\textquotedblright{}.) Say it once more.
\end{tcolorbox}
\section[\texorpdfstring{\ja{はやまる} (hayamaru)}{はやまる (hayamaru)}]{\ja{はやまる} (hayamaru) — to be brought forward}

\label{lesson:JA-C512-hayamaru}



\begin{quote}
Before the new one: say the Japanese for to change, then the Japanese for to switch (over).
\end{quote}

\subsection*{You'll want to know: \ja{はやまる}}
\textbf{\ja{はやまる}} — \emph{hayamaru} — \textquotedblleft{}to be brought forward\textquotedblright{}.

\begin{quote}
\textbf{\ja{もし} \ja{しあいが} \ja{はやまったら、しらせて} \ja{ください。}} — \emph{moshi shiai ga hayamattara, shirasete kudasai} — \textquotedblleft{}If the match is brought forward, please let me know.\textquotedblright{}
\end{quote}

Polite \textbf{\ja{はやまります}}, \emph{hayamarimasu}; a godan verb. Past (polite): \textbf{\ja{かいぎが} \ja{はやまりました。}}, \emph{kaigi ga hayamarimashita}.

\begin{cousinweb}
\textbf{\ja{はやまる}} (\emph{hayamaru}) is a native Japanese word (wago), inherited rather than borrowed from Chinese; its kanji was attached later as a kun'yomi, and its older origin is not settled (JMdict, KANJIDIC2).
\end{cousinweb}

\begin{grammarlens}[title={What you've built}]
One more word for saying what happens if.
\end{grammarlens}

\subsection*{Your turn}
\begin{itemize}
\raggedright
  \item \emph{Say it:} \emph{hayamaru}
  \item \emph{Say it:} \emph{hayamaru}, once more
  \item \emph{Say it:} \emph{moshi shiai ga hayamattara, shirasete kudasai} — If the match is brought forward, please let me know.
  \item \emph{Recall:} say the Japanese for to change, then the Japanese for to switch (over), then say \emph{hayamaru} again
\end{itemize}

\begin{tcolorbox}[breakable,colback=teal!4,colframe=teal!35!black,title={Before you move on}]
What does \ja{はやまる} mean? (\textquotedblleft{}To be brought forward\textquotedblright{}.) Say it once more.
\end{tcolorbox}
\section[\texorpdfstring{\ja{かさなる} (kasanaru)}{かさなる (kasanaru)}]{\ja{かさなる} (kasanaru) — to clash (of plans); to overlap}

\label{lesson:JA-C512-kasanaru}



\begin{quote}
Before the new one: say the Japanese for to switch (over), then the Japanese for to be brought forward.
\end{quote}

\subsection*{You'll want to know: \ja{かさなる}}
\textbf{\ja{かさなる}} — \emph{kasanaru} — \textquotedblleft{}to clash (of plans); to overlap\textquotedblright{}.

\begin{quote}
\textbf{\ja{もし} \ja{よていが} \ja{かさなったら、どちらを} \ja{えらびますか。}} — \emph{moshi yotei ga kasanattara, dochira o erabimasu ka} — \textquotedblleft{}If your plans clash, which will you choose?\textquotedblright{}
\end{quote}

Polite \textbf{\ja{かさなります}}, \emph{kasanarimasu}; a godan verb. Past (polite): \textbf{\ja{しけんと} \ja{しあいが} \ja{かさなりました。}}, \emph{shiken to shiai ga kasanarimashita}.

\begin{cousinweb}
\textbf{\ja{かさなる}} (\emph{kasanaru}) is a native Japanese word (wago), inherited rather than borrowed from Chinese; its kanji was attached later as a kun'yomi, and its older origin is not settled (JMdict, KANJIDIC2).
\end{cousinweb}

\begin{grammarlens}[title={What you've built}]
One more word for saying what happens if.
\end{grammarlens}

\subsection*{Your turn}
\begin{itemize}
\raggedright
  \item \emph{Say it:} \emph{kasanaru}
  \item \emph{Say it:} \emph{kasanaru}, once more
  \item \emph{Say it:} \emph{moshi yotei ga kasanattara, dochira o erabimasu ka} — If your plans clash, which will you choose?
  \item \emph{Recall:} say the Japanese for to switch (over), then the Japanese for to be brought forward, then say \emph{kasanaru} again
\end{itemize}

\begin{tcolorbox}[breakable,colback=teal!4,colframe=teal!35!black,title={Before you move on}]
What does \ja{かさなる} mean? (\textquotedblleft{}To clash (of plans); to overlap\textquotedblright{}.) Say it once more.
\end{tcolorbox}
\section[\texorpdfstring{\ja{ひがわり} (higawari)}{ひがわり (higawari)}]{\ja{ひがわり} (higawari) — changing every day; of the day}

\label{lesson:JA-C512-higawari}



\begin{quote}
Before the new one: say the Japanese for to be brought forward, then the Japanese for to clash (of plans).
\end{quote}

\subsection*{You'll want to know: \ja{ひがわり}}
\textbf{\ja{ひがわり}} — \emph{higawari} — \textquotedblleft{}changing every day; of the day\textquotedblright{}.

\begin{quote}
\textbf{\ja{ねだんは} \ja{ひがわりの} \ja{りょうりに} \ja{よります。}} — \emph{nedan wa higawari no ryōri ni yorimasu} — \textquotedblleft{}The price depends on the dish of the day.\textquotedblright{}
\end{quote}

\begin{cousinweb}
\textbf{\ja{ひがわり}} (\emph{higawari}) is a native Japanese compound (wago), inherited rather than borrowed from Chinese: built from \emph{hi}, \textquotedblleft{}day\textquotedblright{} and \emph{gawari}, \textquotedblleft{}changing\textquotedblright{} (JMdict, KANJIDIC2).
\end{cousinweb}

\begin{grammarlens}[title={What you've built}]
One more word for saying what happens if.
\end{grammarlens}

\subsection*{Your turn}
\begin{itemize}
\raggedright
  \item \emph{Say it:} \emph{higawari}
  \item \emph{Say it:} \emph{higawari}, once more
  \item \emph{Say it:} \emph{nedan wa higawari no ryōri ni yorimasu} — The price depends on the dish of the day.
  \item \emph{Recall:} say the Japanese for to be brought forward, then the Japanese for to clash (of plans), then say \emph{higawari} again
\end{itemize}

\begin{tcolorbox}[breakable,colback=teal!4,colframe=teal!35!black,title={Before you move on}]
What does \ja{ひがわり} mean? (\textquotedblleft{}Changing every day; of the day\textquotedblright{}.) Say it once more.
\end{tcolorbox}
